% 替换位置：第9章“这里是模型评价”占位句；以下为与本研究结果有关的评价段落
% 这是修改片段，不是完整文档；不要整份直接追加到原文末尾。
本模型统一四问的地形、载荷、能耗和资源占用规则，并用独立回放检查计算结果，
有助于避免逐段载荷、满电复用和跨问继承口径不一致。
批次重组、机型选择和通信服务日历共同参与求解，
第二问通过24个运输架次完成80箱按期交付，
第三问配置3个中继架次完成全过程数值通信保障。

模型也存在限制。第二、三问尚未穷尽货箱分配、运输顺序和连续中继位置，
零延误仅证明第一层目标达到下界，不能证明完整字典序全局最优。
通信检查仍为有限分辨率的数值验证，不能严格排除采样间隙失效。
能量标定、爬升效率、未计运输交接悬停耗能及无线并发容量等工作假设需要进一步核验。

此外，分区要求资源不跨组调配，因而会产生额外配置及中继任务复制能耗。
第四问的最优性仅针对冻结时序、完整中继复制和组内重新编号口径，
不能扩展为运输、通信与分区共同决策时的全局最优性。
